\section{Segment Tree}

\textbf{Bai toan.}
Cho mang $a[0 \ldots n-1]$ gom $n$ so nguyen. Ho tro hai thao tac:
\begin{itemize}
  \item \texttt{update(i, v)}: gan $a[i] = v$.
  \item \texttt{query(l, r)}: tinh $\displaystyle\sum_{k=l}^{r} a[k]$.
\end{itemize}
Yeu cau: moi thao tac phai chay trong $O(\log n)$.

\subsection{Cau truc cay}

Segment tree la cay nhi phan day du, trong do moi node dai dien cho mot doan $[lo, hi]$ cua mang.
Node la luu $a[i]$, con node trong luu tong cua hai con:

\[
  \text{node}[k] = \text{node}[2k] + \text{node}[2k+1]
\]

Voi $n = 6$ va $a = [1, 3, 5, 7, 9, 2]$, cay co dang:

\begin{diagram}[id="segtree-static"]
\shape{T}{Tree}{
  data=[1,3,5,7,9,2],
  kind="segtree",
  show_sum=true,
  label="Segment tree cho a = [1,3,5,7,9,2]"
}
\recolor{T.node["[0,5]"]}{state=good}
\recolor{T.node["[0,2]"]}{state=idle}
\recolor{T.node["[3,5]"]}{state=idle}
\end{diagram}

\subsection{Phan tich do phuc tap}

Xay cay (\texttt{build}): duyet toan bo $2n - 1$ node, moi node thuc hien phep cong hang so
$\Rightarrow O(n)$.

Truy van tong (\texttt{query}): tai moi tang cua cay, ta chi di vao toi da $2$ nhanh.
Cay co chieu cao $\lceil \log_2 n \rceil$, nen moi truy van ton $O(\log n)$.

Cap nhat diem (\texttt{update}): di tu la len root, cap nhat tong doc duong di.
Duong di co do dai $O(\log n)$.

\subsection{Minh hoa: truy van sum(1, 4)}

\begin{animation}[id="segtree-query", label="Segment Tree -- Range sum query(1,4)"]
\shape{a}{Array}{size=6, data=[1,3,5,7,9,2], labels="0..5", label="$a$"}
\shape{T}{Tree}{
  data=[1,3,5,7,9,2],
  kind="segtree",
  show_sum=true,
  label="Segment tree"
}

\step
\narrate{Mang $a = [1, 3, 5, 7, 9, 2]$. Segment tree da duoc xay: moi node luu tong doan ma no dai dien. Root \texttt{[0,5]} = 27. Ta can tinh $\text{sum}(1, 4) = a[1] + a[2] + a[3] + a[4] = 3 + 5 + 7 + 9 = 24$.}

\step
\recolor{T.node["[0,5]"]}{state=current}
\narrate{Bat dau tu root \texttt{[0,5]}. Doan $[0,5]$ khong nam hoan toan trong $[1,4]$, nen ta phai di xuong ca hai nhanh con.}

\step
\recolor{T.node["[0,5]"]}{state=dim}
\recolor{T.node["[0,2]"]}{state=current}
\narrate{Di xuong nhanh trai: node \texttt{[0,2]}. Doan $[0,2]$ khong nam hoan toan trong $[1,4]$ (vi co phan tu $0$), nen tiep tuc chia.}

\step
\recolor{T.node["[0,2]"]}{state=dim}
\recolor{T.node["[0,0]"]}{state=current}
\narrate{Di xuong nhanh trai cua \texttt{[0,2]}: node \texttt{[0,0]}. Doan $[0,0]$ nam hoan toan ngoai $[1,4]$ -- bo qua node nay, tra ve $0$.}

\step
\recolor{T.node["[0,0]"]}{state=dim}
\recolor{T.node["[1,2]"]}{state=current}
\narrate{Di xuong nhanh phai cua \texttt{[0,2]}: node \texttt{[1,2]}. Doan $[1,2]$ nam hoan toan trong $[1,4]$ -- lay ngay gia tri $3 + 5 = 8$. Khong can di xuong sau hon.}

\step
\recolor{T.node["[1,2]"]}{state=good}
\recolor{T.edge[("[0,2]","[1,2]")]}{state=good}
\narrate{Node \texttt{[1,2]} dong gop $8$ vao ket qua. Danh dau \texttt{good}. Quay lai root de di sang nhanh phai.}

\step
\recolor{T.node["[3,5]"]}{state=current}
\narrate{Di xuong nhanh phai cua root: node \texttt{[3,5]}. Doan $[3,5]$ khong nam hoan toan trong $[1,4]$ (vi co phan tu $5$), nen tiep tuc chia.}

\step
\recolor{T.node["[3,5]"]}{state=dim}
\recolor{T.node["[3,4]"]}{state=current}
\narrate{Di xuong nhanh trai cua \texttt{[3,5]}: node \texttt{[3,4]}. Doan $[3,4]$ nam hoan toan trong $[1,4]$ -- lay ngay gia tri $7 + 9 = 16$.}

\step
\recolor{T.node["[3,4]"]}{state=good}
\recolor{T.edge[("[3,5]","[3,4]")]}{state=good}
\recolor{T.node["[5,5]"]}{state=current}
\narrate{Node \texttt{[3,4]} dong gop $16$. Di sang nhanh phai cua \texttt{[3,5]}: node \texttt{[5,5]}. Doan $[5,5]$ nam hoan toan ngoai $[1,4]$ -- bo qua, tra ve $0$.}

\step
\recolor{T.node["[5,5]"]}{state=dim}
\recolor{T.node["[0,5]"]}{state=done}
\recolor{T.node["[1,2]"]}{state=good}
\recolor{T.node["[3,4]"]}{state=good}
\recolor{T.node["[0,0]"]}{state=dim}
\recolor{T.node["[0,2]"]}{state=dim}
\recolor{T.node["[3,5]"]}{state=dim}
\recolor{T.node["[5,5]"]}{state=dim}
\narrate{Ket qua: $\text{sum}(1,4) = 8 + 16 = 24$. Chi co $2$ node (\texttt{[1,2]} va \texttt{[3,4]}, to \texttt{good}) thuc su dong gop vao ket qua. Cac node con lai hoac bi bo qua hoac chi la trung gian. Tong so node tham: $O(\log n)$.}
\end{animation}

\subsection{Cap nhat diem}

Khi thay doi $a[i] = v$, ta chi can cap nhat node la tuong ung va truyen nguoc len root.
Cu the, tu la $[i, i]$, di len cha $[lo, hi]$, tinh lai:
\[
  \text{node}[\text{cha}] = \text{node}[\text{con\_trai}] + \text{node}[\text{con\_phai}]
\]
Do cay co chieu cao $O(\log n)$, thao tac update cung ton $O(\log n)$.

\subsection{Cai dat}

\begin{lstlisting}[language=C++]
const int MAXN = 200005;
int n, a[MAXN], tree[4 * MAXN];

void build(int node, int lo, int hi) {
    if (lo == hi) {
        tree[node] = a[lo];
        return;
    }
    int mid = (lo + hi) / 2;
    build(2 * node, lo, mid);
    build(2 * node + 1, mid + 1, hi);
    tree[node] = tree[2 * node] + tree[2 * node + 1];
}

int query(int node, int lo, int hi, int l, int r) {
    if (r < lo || hi < l) return 0;       // hoan toan ngoai
    if (l <= lo && hi <= r) return tree[node]; // hoan toan trong
    int mid = (lo + hi) / 2;
    return query(2 * node, lo, mid, l, r)
         + query(2 * node + 1, mid + 1, hi, l, r);
}

void update(int node, int lo, int hi, int pos, int val) {
    if (lo == hi) {
        tree[node] = val;
        return;
    }
    int mid = (lo + hi) / 2;
    if (pos <= mid)
        update(2 * node, lo, mid, pos, val);
    else
        update(2 * node + 1, mid + 1, hi, pos, val);
    tree[node] = tree[2 * node] + tree[2 * node + 1];
}

// Su dung:
// build(1, 0, n - 1);
// query(1, 0, n - 1, l, r);
// update(1, 0, n - 1, i, v);
\end{lstlisting}

\textbf{Luu y:}
Kich thuoc mang \texttt{tree} can $4n$ phan tu (khong phai $2n$) do cach danh so $1$-indexed
va padding khi $n$ khong phai luy thua cua $2$.
Do phuc tap bo nho: $O(n)$.
Do phuc tap tong the: build $O(n)$, query $O(\log n)$, update $O(\log n)$.
